\section*{Bab \ref{ch:b3:milky-way-galaxies} --- Bima Sakti dan Galaksi-Galaksi}

\begin{solution}{exo:b3:milky-way-galaxies:1}
(a)~$2\pi R_0 = 2\pi\times\qty{2.523e17}{km} = \qty{1.585e18}{km}$; pada
\qty{229}{km.s^{-1}} ini membutuhkan \qty{6.92e15}{s}, yaitu
\qty{219}{Ma}. (b)~$4540/219 = 20.7$: sekitar dua puluh putaran.
\end{solution}

\begin{solution}{exo:b3:milky-way-galaxies:2}
(a)~$\qty{1}{au}/\qty{1}{yr} = \num{1.496e8}/\num{3.156e7} =
\qty{4.74}{km.s^{-1}}$. (b)~$v = 4.74\times6.411\times8.178 =
\qty{248.5}{km.s^{-1}}$, berada dalam rentang tinjauan
\qty[separate-uncertainty=true]{248 \pm 3}{km.s^{-1}}. (c)~$248.5 - 11 = \qty{237.5}{km.s^{-1}}$,
sekitar \qty{8}{km.s^{-1}} di atas \qty{229}{km.s^{-1}} milik raksasa merah:
keduanya bertumpu pada nilai $R_0$ yang berbeda (8.178 dan \qty{8.122}{kpc}),
dan \omterm{def:b3:milky-way-galaxies:lsr}{standar diam lokal} sendiri dapat bergeser hingga sekitar \qty{15}{km.s^{-1}} relatif
terhadap orbit yang benar-benar melingkar.
\end{solution}

\begin{solution}{exo:b3:milky-way-galaxies:3}
(a)~$n = 10(1 - 0.6) = 4$: E4. (b)~$b/a = 1 - 7/10 = 0.3$. (c)~Kelas
memakai rasio sumbu tampak: elipsoid pipih yang dilihat sepanjang sumbu pendeknya
memiliki kontur melingkar dan digolongkan sebagai E0.
\end{solution}

\begin{solution}{exo:b3:milky-way-galaxies:4}
$d = 10^{\mu/5 + 1}$~pc: LMC \qty{49.6}{kpc}, Andromeda
\qty{783}{kpc}, NGC~4258 \qty{7.58}{Mpc}. Sebuah bintang di Awan memiliki
paralaks $1/\qty{49.6}{kpc} = \qty{0.020}{mas}$, setara dengan
ketidakpastian katalog terbaik untuk bintang terang: satu paralaks saja tidak
memberikan jarak sama sekali, dan Awan ditentukan dengan cara lain
(binari gerhana), dari situlah skala Cepheid dimulai.
\end{solution}

\begin{solution}{exo:b3:milky-way-galaxies:5}
(a)~$\omega_0 = A - B = \qty{27.2}{km.s^{-1}.kpc^{-1}}$ dan $V_0 =
27.2\times8.178 = \qty{222}{km.s^{-1}}$. (b)~$V'_0 = -(A + B) =
\qty{-3.4}{km.s^{-1}.kpc^{-1}}$: kurva yang turun perlahan. (c)~Keduanya
berada di bawah nilai raksasa merah (229 dan \qty{-1.7}{km.s^{-1}} per
kiloparsek); para penulis kurva itu memperoleh hal yang sama, dengan $V_0$ sekitar dua
simpangan baku lebih rendah, yang diturunkan dari $R_0$ dan \omterm{def:b3:milky-way-galaxies:rotation-curve}{konstanta Oort} ini. Rumus Oort
mengandaikan orbit melingkar dalam cakram simetri-aksial; $C$ dan
$K$ yang tidak nol menunjukkan bahwa medan kecepatan lokal terganggu oleh batang dan
lengan spiral, dan $A$ serta $B$ lokal membawa gangguan tersebut.
\end{solution}

\begin{solution}{exo:b3:milky-way-galaxies:6}
(a)~$\rho = V^2/(4\pi GR^2) = 229^2/(4\pi\times\num{4.301e-6}\times
8.178^2) = \num{1.45e7}$ massa Matahari per kiloparsek kubik, atau
\num{0.0145} per parsek kubik. (b)~$M_\odot c^2 = \num{1.9884e30}\times
(\num{2.998e8})^2 = \qty{1.787e47}{J} = \qty{1.115e57}{GeV}$, dan
$\qty{1}{pc^3} = (\qty{3.0857e18}{cm})^3 = \qty{2.938e55}{cm^3}$: rasionya
adalah \qty{37.96}{GeV.cm^{-3}}. Maka $0.30/37.96 = \num{0.0079}$ massa Matahari
per parsek kubik. (c)~Massa jenis kurva datar menghitung seluruh massa
yang menahan orbit, termasuk bintang dan gas, seolah-olah tersebar dalam
bola; materi gelap hanya sebagian darinya, kira-kira separuh di sini.
\end{solution}

\begin{solution}{exo:b3:milky-way-galaxies:7}
(a)~$W_I(\mathrm{LMC}) = -3.314\times1.4771 + 15.888 = \qty{10.993}{mag}$,
maka $M_W = 10.993 - 18.477 = \qty{-7.484}{mag}$. (b)~$\mu = 25.00 + 7.484 =
32.484$ dan $d = 10^{32.484/5 + 1}~\mathrm{pc} = \qty{31.4}{Mpc}$.
(c)~$\delta d/d = (\ln 10/5)\,\delta\mu = 0.4605\times0.077 =
\qty{3.5}{\percent}$ untuk satu bintang, $3.5/\sqrt{50} = \qty{0.5}{\percent}$
untuk lima puluh bintang, jika kesalahannya saling bebas (kesalahan jarak Awan
itu sendiri ikut menambah semuanya secara sama).
\end{solution}

\begin{solution}{exo:b3:milky-way-galaxies:8}
$\mu = I - A_I - M_I = 26.00 - 0.05 + 4.05 = 30.00$, maka $d =
10^{7}~\mathrm{pc} = \qty{10.0}{Mpc}$. Kesalahan sebesar \qty{0.04}{mag} adalah
$0.4605\times0.04 = \qty{1.8}{\percent}$ pada jarak.
\end{solution}

\begin{solution}{exo:b3:milky-way-galaxies:9}
(a)~$d = cz/H_0$ memberikan $5\log_{10}d = 5\log_{10}(cz) - 5\log_{10}H_0$,
maka $m_B - M_B - 25 = 5\log_{10}(cz) - 5\log_{10}H_0$; dengan
$\log_{10}(cz) = a_B + 0.2\,m_B$, $5\log_{10}H_0 = 5a_B + M_B + 25$, yang
merupakan hasilnya. (b)~$\log_{10}H_0 = -3.8506 + 0.714158 + 5 = 1.8636$, maka
$H_0 = \qty{73.04}{km.s^{-1}.Mpc^{-1}}$. (c)~$M_B = 5(\log_{10}67.4 -
5.714158) = \qty{-19.427}{mag}$: supernova lebih terang sebesar \qty{0.174}{mag},
maka seluruh jarak lebih besar dengan faktor $10^{0.2\times0.174} = 1.084$, yaitu
rasio $73.04/67.4$. Ketegangan ini adalah perbedaan delapan persen dalam
skala jarak.
\end{solution}

\begin{solution}{exo:b3:milky-way-galaxies:10}
(a)~Bintang di tepi, dalam orbit melingkar, memiliki $V^2/R \approx GM/R^2$
(pendekatan bola). (b)~$R = \sqrt{M/(\pi\Sigma)}$, sehingga $V^2 =
GM/R = G\sqrt{\pi\Sigma}\,M^{1/2}$ dan $M = V^4/(\pi\Sigma G^2) \propto
V^4$; dengan rasio massa-terang yang sama, $L \propto V^4$ juga.
(c)~$(100/229)^4 = \num{0.036}$: galaksi yang lebih cepat 27 kali lebih
masif. (d)~Luminositas sama: fluks turun sebanding $1/d^2$, sehingga yang lebih redup
berada pada jarak dua kali lipat.
\end{solution}

\begin{solution}{exo:b3:milky-way-galaxies:11}
(a)~$\qty{1}{km.s^{-1}} = \qty{1.0227}{kpc.Ga^{-1}}$, maka
\qty{229}{km.s^{-1}} sama dengan \qty{234.2}{kpc.Ga^{-1}}. (b)~$2\pi R/V$:
\qty{107}{Ma} pada \qty{4}{kpc}, \qty{322}{Ma} pada \qty{12}{kpc}.
(c)~$\Delta\phi = Vt\,(1/4 - 1/12) = 234.2\times0.1667 = \qty{39.0}{rad}$,
6.2 putaran: lengan ini akan menjadi gulungan enam lilitan rapat antara 4 dan
\qty{12}{kpc}, jauh lebih rapat daripada lengan yang teramati, dan semakin rapat lagi
pada setiap rotasi berikutnya. (d)~Dalam gelombang massa jenis, pola itu berputar dengan satu
laju sudut, seperti benda tegar; bintang dan gas melewatinya, sehingga
bentuk lengan tidak terlilit.
\end{solution}

\begin{solution}{exo:b3:milky-way-galaxies:12}
(a)~$a = (4.2\times10^6\times16^2)^{1/3} = \qty{1024}{au}$; pada
\qty{8178}{pc} membentuk sudut $1024/8178 = \ang{;;0.125}$ (satu au pada satu parsek
berukuran satu detik busur). (b)~$R_S = 2GM/c^2 = 2\times\num{6.674e-11}
\times\num{8.35e36}/(\num{2.998e8})^2 = \qty{1.24e10}{m} = \qty{0.083}{au}$;
diameter $2R_S$ membentuk sudut $0.166/8178 = \num{2.0e-5}$ detik busur, sekitar
\qty{20}{\micro as}. (c)~$R_S = \qty{4.02e-7}{pc}$, dan $M/(\tfrac43\pi
R_S^3) = \num{1.5e25}$ massa Matahari per parsek kubik, empat puluh kali
batas bawah yang ditetapkan oleh gerak dan ukuran Sgr~A*: batas itu berada dalam
faktor empat puluh dari massa jenis lubang hitam itu sendiri.
\end{solution}

\begin{solution}{pb:b3:milky-way-galaxies:1}

\textbf{Bagian I.}

\textbf{1.} $v_r = A\,d\sin2\ell$: $+A d$ pada \ang{45}, 0 pada \ang{90}, $-A d$ pada
\ang{135}, 0 pada \ang{180}.

\textbf{2.} $A = (15.4 + 15.2)/(2\times\qty{1.00}{kpc}) =
\qty{15.3}{km.s^{-1}.kpc^{-1}}$.

\textbf{3.} Pada \ang{90}, $v_t = d(B - A) = 4.74\times(-5.91)\times1.00 =
\qty{-28.02}{km.s^{-1}}$, sehingga $B = 15.3 - 28.02 =
\qty{-12.72}{km.s^{-1}.kpc^{-1}}$.

\textbf{4.} Pada \ang{0}, $v_t = d(A + B) = 4.74\times0.55 =
\qty{2.61}{km.s^{-1}}$, sehingga $B = 2.61 - 15.3 = \num{-12.69}$: keduanya
sepakat dalam \num{0.03}.

\textbf{5.} $\omega_0 = A - B = \qty{28.0}{km.s^{-1}.kpc^{-1}}$ dan
$V_0 = 28.0\times8.18 = \qty{229}{km.s^{-1}}$.

\textbf{6.} $V'_0 = -(A + B) = \qty{-2.6}{km.s^{-1}.kpc^{-1}}$: menurun perlahan.

\textbf{7.} $P = 2\pi/\omega_0$: $2\pi\times\qty{8.18}{kpc}/\qty{229.2}{km.s^{-1}}
= \qty{219}{Ma}$; $4540/219 = 20.7$ putaran.

\textbf{Bagian II.}

\textbf{8.} Titik singgung memerlukan garis pandang yang melewati lebih dekat ke
pusat daripada Matahari, yaitu $\sin\ell > 0$ dengan $\ell < \ang{90}$ (atau
kuadran cerminnya, $\ang{270} < \ell < \ang{360}$, tempat kecepatan terbesar
paling negatif). Di luar lingkaran Matahari, jarak ke pusat bertambah sepanjang
garis pandang: tidak ada titik yang menonjol, dan kecepatan saja tidak
menunjukkan di mana gas berada.

\textbf{9.} $R = 8.18\sin\ell$: \qtylist{5.26;6.27;7.08;7.69}{kpc}.

\textbf{10.} $V = v_\mathrm{max} + V_0\sin\ell$: $82.9 + 147.2 = 230.1$,
$55.6 + 175.4 = 231.0$, $32.4 + 198.3 = 230.7$, $14.9 + 215.2 =
\qty{230.1}{km.s^{-1}}$.

\textbf{11.} $\Delta\nu = \nu v/c = \num{1.420406e9}\times82.9/299792 =
\qty{393}{kHz}$, ke arah frekuensi lebih rendah: gas menjauh.

\textbf{12.} Dalam sekitar \qty{3}{km.s^{-1}} dari raksasa merah: antara 5 dan
\qty{8}{kpc} kurvanya datar.

\textbf{13.} Massa titik memberikan $V = V_0\sqrt{R_0/R} = 229\times\sqrt{8.18/5.26} =
\qty{286}{km.s^{-1}}$ pada titik singgung pertama, jauh lebih cepat daripada
\qty{230}{km.s^{-1}} yang ditemukan: massanya tidak terpusat di pusat,
sebagian besar berada antara 5 dan \qty{8}{kpc} serta di luarnya.

\textbf{Bagian III.}

\textbf{14.} $\num{1.067e10} + 2\times\num{3.944e10} = \num{8.96e10}$ massa
Matahari, hampir dua kali lipat \num[separate-uncertainty=true]{5 \pm 1 e10} dari ulasan; bintang-bintang
model mungkin terlalu berat, yang akan membuat materi gelap yang ditemukan di bawah
menjadi taksiran yang terlalu rendah.

\textbf{15.} $V_b^2 = GMR^2/(R^2 + b^2)^{3/2} = \num{4.301e-6}\times\num{1.067e10}
\times400/(400.09)^{3/2}$, $V_b = \qty{47.9}{km.s^{-1}}$.

\textbf{16.} Tipis: $(a + b)^2 = 30.80$, $V^2 = \num{4.301e-6}\times\num{3.944e10}
\times400/(430.8)^{3/2}$, $V = \qty{87.1}{km.s^{-1}}$. Tebal:
$(a + b)^2 = 11.56$, $V = \qty{90.1}{km.s^{-1}}$.

\textbf{17.} $V_\mathrm{vis} = \sqrt{47.9^2 + 87.1^2 + 90.1^2} =
\qty{134.2}{km.s^{-1}}$.

\textbf{18.} Pada \qty{8.18}{kpc}: 74.8, 108.4 dan \qty{127.8}{km.s^{-1}},
bersama-sama \qty{183.5}{km.s^{-1}}, dibandingkan dengan $V_0 = \qty{229}{km.s^{-1}}$:
massa tampak sudah kurang di Matahari.

\textbf{19.} $\sqrt{G\times\num{8.96e10}/20} = \qty{138.8}{km.s^{-1}}$. Nilai
model \qty{134.2}{km.s^{-1}} cukup dekat, karena pada \qty{20}{kpc},
empat sampai enam panjang skala cakram di luarnya, hampir seluruh massa tampak
berada di dalam orbit: kurva tampak kini berbentuk Keplerian.

\textbf{Bagian IV.}

\textbf{20.} $M = 20\times200^2/\num{4.301e-6} = \num{1.860e11}$ massa Matahari.

\textbf{21.} $M_\mathrm{vis} = 20\times134.2^2/\num{4.301e-6} = \num{8.37e10}$
massa Matahari.

\textbf{22.} $M_\mathrm{dark} = \num{1.860e11} - \num{8.37e10} = \num{1.02e11}$;
karena kedua massa adalah $RV^2/G$ pada $R$ yang sama, pecahannya adalah
$1 - (V_\mathrm{vis}/V)^2 = 1 - (134.2/200)^2 = \num{0.550}$.

\textbf{23.} $x = 20/14.8 = 1.351$, $\ln(2.351) - 1.351/2.351 = 0.8549 - 0.5748
= 0.2801$, $M = 4\pi\times\num{1.06e7}\times14.8^3\times0.2801 =
\num{1.21e11}$ massa Matahari; pecahan materi gelapnya adalah
$1.21/(1.21 + 0.837) = \num{0.59}$. Halo dicocokkan dengan seluruh
kurva, yang pada \qty{20}{kpc} ternyata melampaui nilai itu: pengukuran di sana
turun sekitar \qty{10}{km.s^{-1}} di bawah model halus (\qty{209}{km.s^{-1}}).

\textbf{24.} $x = 8.18/14.8 = 0.553$, $\rho = \num{1.06e7}/(0.553\times1.553^2)
= \num{7.95e6}$ massa Matahari per kiloparsek kubik, \num{0.0080} per
parsek kubik, atau \qty{0.30}{GeV.cm^{-3}}, nilai yang dikutip oleh para penulis
pencocokan itu.

\textbf{25.} $1 - 8.37/18.60 = \num{0.550}$: pecahan materi gelap dari massa
dalam \qty{20}{kpc} adalah
\textbf{\qty[text-series-to-math=true]{55}{\percent}}.
\end{solution}
